\setcounter{section}{21}

  \zwischenueberschrift{Soal latihan}

\inputaufgabe
{}
{

Deskripsikan berbagai jenis pakaian sebagai
\definitionsverweis {manifold berbatas}{}{}
berdimensi dua.

}
{} {}

\inputaufgabe
{}
{

Definisikan konsep pemetaan diferensiabel, difeomorfisme, ruang singgung, bundel tangen, dan seterusnya untuk
\definitionsverweis {manifold berbatas}{}{.}

}
{} {}

\inputaufgabe
{}
{

Deskripsikan
\definitionsverweis {interval riil}{}{}
sebagai
\definitionsverweis {manifold berbatas}{}{.}
Apa
\definitionsverweis {batas}{}{}
masing-masing interval tersebut, dan interval mana saja yang saling
\definitionsverweis {difeomorfik}{}{?}

}
{} {}

\inputaufgabe
{}
{

$\,$

\bild{ \begin{center}
\includegraphics[width=5.5cm]{ \bildeinlesungsvg {Circle_on_sphere_wireframe_10deg_6r} {svg} }
\end{center}
\bildtext {} }

\bildlizenz { Circle on sphere wireframe 10deg 6r.svg } {} {Itai} {Commons} {CC BY 3.0} {}

Tunjukkan bahwa potongan permukaan berwarna biru maupun merah, masing-masing bersama garis pembatasnya, merupakan
\definitionsverweis {manifold berbatas}{}{.}
Apa batasnya? Apakah kedua manifold tersebut
\definitionsverweis {difeomorfik}{}{?}
Adakah manifold yang lebih sederhana dan difeomorfik dengannya?

}
{} {}

\inputaufgabe
{}
{

Misalkan $M$ suatu
\definitionsverweis {manifold diferensiabel}{}{}
\zusatzklammer {tanpa batas} {} {}
dan $N$ suatu
\definitionsverweis {manifold diferensiabel berbatas}{}{.}
Apa yang dapat dikatakan tentang
\definitionsverweis {produk}{}{}
\mathl{M \times N}{?}

}
{} {}

\inputaufgabe
{}
{

Tunjukkan bahwa anulus tertutup

\mathdisp {\left\{ P \in \R^2  \mid   1 \leq \Vert { P } \Vert \leq 2         \right\}} {  }

dan silinder tertutup

\mathdisp {S^1 \times [0,1]} {  }

merupakan
\definitionsverweis {manifold berbatas}{}{}
yang saling
\definitionsverweis {difeomorfik}{}{.}

}
{} {}

\inputaufgabegibtloesung
{}
{

Jelaskan mengapa $\R^2$ tidak dapat diberi struktur
\definitionsverweis {manifold berbatas}{}{}
sedemikian rupa sehingga himpunan bagian $T$ yang diberikan menjadi
\definitionsverweis {batas}{}{}
dari struktur tersebut.
\aufzaehlungvierabc{
\mavergleichskette
{\vergleichskette
{ T
}
{ = }{ {]0,1[}
}
{  }{
}
{    }{
}
{   }{
}
}
{}{}{,}
dengan interval tersebut terletak pada sumbu $x$.
}{
\mavergleichskette
{\vergleichskette
{ T
}
{ = }{ [0,1]
}
{  }{
}
{    }{
}
{   }{
}
}
{}{}{,}
dengan interval tersebut terletak pada sumbu $x$.
}{
\mavergleichskette
{\vergleichskette
{ T
}
{ = }{ \R
}
{  }{
}
{    }{
}
{   }{
}
}
{}{}{,}
dengan $\R$ menyatakan sumbu $x$.
}{
\mavergleichskette
{\vergleichskette
{ T
}
{ = }{ S^1
}
{  }{
}
{    }{
}
{   }{
}
}
{}{}{.}
}

}
{} {}

\inputaufgabe
{}
{

Tunjukkan bahwa kedua
\definitionsverweis {difeomorfisme}{}{}
yang saling invers,
\maabbeledisp {\varphi} { {\mathbb H} } { U \left(   0 , 1  \right)
} { z} { { \frac{  z-  { \mathrm i}  }{ z + { \mathrm i}  } }
} {,}
dan
\maabbeledisp {\psi} { U \left(   0 , 1  \right)  } { {\mathbb H}
} { w } { { \frac{   (w+1)  { \mathrm i}  }{  1-w  } }
} {,}
\zusatzklammer {lihat
[[Obere Halbebene/Einheitskreis/Rational isomorph/Fakt/Beweis/Aufgabe|Soal 18.10]]} {} {,}
dapat diperluas menjadi difeomorfisme
\zusatzklammer {antara
\definitionsverweis {manifold berbatas}{}{}} {} {}
antara setengah bidang tertutup dan cakram tertutup berlubang
\mathl{B  \left(  0 , 1 \right) \setminus \{  \left(  1   , \,   0                   \right) \}}{}.

}
{} {}

\inputaufgabe
{}
{

Misalkan $M$ suatu
\definitionsverweis {manifold diferensiabel berbatas}{}{.}
Yang dimaksud dengan \stichwort {lintasan-separuh diferensiabel} {} ialah setiap
\definitionsverweis {pemetaan diferensiabel}{}{}
\maabbdisp {\gamma} { [0, \epsilon [ } { M
} {}
atau
\maabbdisp {\gamma} {] -\epsilon, 0] } { M
} {}
\zusatzklammer {dengan \mathlk{\epsilon >0}{;} masing-masing disebut mengarah ke dalam dan mengarah ke luar} {} {.}
Definisikan kapan dua lintasan-separuh dengan
\mathl{\gamma_1(0) = \gamma_2(0) =P \in M}{} disebut \stichwort {ekuivalen secara tangensial} {,} lalu tunjukkan bahwa definisi ini memberikan suatu relasi ekuivalensi dan bahwa
\definitionsverweis {himpunan hasil bagi}{}{}
yang diperoleh merupakan ruang vektor riil yang disebut \stichwort {ruang singgung} {} di $P$. Karakterisasikan kelas-kelas ekuivalensi yang dapat direpresentasikan baik oleh lintasan-separuh yang mengarah ke dalam maupun yang mengarah ke luar.

}
{} {}

\inputaufgabegibtloesung
{}
{

Misalkan $M$ suatu
\definitionsverweis {manifold berbatas}{}{}
dan

\mavergleichskette
{\vergleichskette
{ P
}
{ \in }{ \partial M
}
{  }{
}
{    }{
}
{   }{
}
}
{}{}{}
suatu titik batas. Tunjukkan bahwa untuk suatu vektor singgung

\mavergleichskette
{\vergleichskette
{ v
}
{ \in }{ T_PM
}
{  }{
}
{    }{
}
{   }{
}
}
{}{}{}
kedua sifat berikut ekuivalen.
\aufzaehlungzwei {Terdapat lintasan yang terdiferensialkan secara kontinu
\maabb {\gamma} { [- \epsilon, 0]} { M
} {}
dengan

\mavergleichskette
{\vergleichskette
{ \gamma(0)
}
{ = }{ P
}
{  }{
}
{    }{
}
{   }{
}
}
{}{}{}
dan

\mavergleichskette
{\vergleichskette
{ \gamma'(0)
}
{ = }{ v
}
{  }{
}
{    }{
}
{   }{
}
}
{}{}{.}
} {Pada setiap peta dengan setengah ruang negatif, vektor $v$ dipetakan oleh pemetaan singgung ke setengah ruang kanan.
}

}
{} {}

\inputaufgabe
{}
{

Misalkan $M$ suatu
\definitionsverweis {manifold diferensiabel berbatas}{}{}
dan
\maabbdisp {\gamma} {{]{-\epsilon}, \epsilon[}} {M
} {}
suatu
\definitionsverweis {lintasan diferensiabel}{}{}
dengan

\mavergleichskettedisp
{\vergleichskette
{ \gamma (0)
}
{ =} { P
}
{ \in} { \partial M
}
{ } {
}
{ } {
}
}
{}{}{.}
Tunjukkan bahwa
\mathl{\gamma'(0) \in T_P \partial M}{.}

}
{} {}

\inputaufgabegibtloesung
{}
{

Misalkan $M$ suatu
\definitionsverweis {manifold berbatas}{}{.}
Tunjukkan bahwa setiap himpunan bagian terbuka

\mavergleichskette
{\vergleichskette
{ N
}
{ \subseteq }{ M
}
{  }{
}
{    }{
}
{   }{
}
}
{}{}{}
juga merupakan manifold dengan batas
\zusatzklammer {yang mungkin kosong} {} {,}
dan bahwa

\mavergleichskettedisp
{\vergleichskette
{ \partial N
}
{ =} { \partial M \cap N
}
{ } {
}
{ } {
}
{ } {
}
}
{}{}{}
berlaku.

}
{} {}

\inputaufgabegibtloesung
{}
{

Misalkan

\mavergleichskette
{\vergleichskette
{ H
}
{ \subseteq }{ \R^n
}
{  }{
}
{    }{
}
{   }{
}
}
{}{}{}
suatu
\definitionsverweis {setengah ruang Euklides}{}{}
dan

\mavergleichskette
{\vergleichskette
{ P
}
{ \in }{ H
}
{  }{
}
{    }{
}
{   }{
}
}
{}{}{.}
Misalkan terdapat himpunan yang terbuka di $\R^n$, yaitu $V$, dengan

\mavergleichskette
{\vergleichskette
{ P
}
{ \in }{ V
}
{ \subseteq }{ H
}
{    }{
}
{   }{
}
}
{}{}{.}
Tunjukkan bahwa
\mathl{P}{} bukan titik batas $H$.

}
{} {}

\inputaufgabe
{}
{

Misalkan
\mathl{U_1 \subseteq H_1 \subset \R^n}{} dan
\mathl{U_2 \subseteq H_2 \subset \R^n}{}
merupakan
\definitionsverweis {himpunan bagian terbuka}{}{}
di dalam
\definitionsverweis {setengah ruang Euklides}{}{}
\mathkor {} {H_1} {dan} {H_2} {,}
serta
\maabbdisp {\varphi} {U_1 } { U_2
} {}
suatu
\definitionsverweis {difeomorfisme}{}{.}
Tunjukkan bahwa untuk setiap titik
\mathl{P\in U_1}{} terdapat lingkungan terbuka
\mathl{P\in V_1 \subset \R^n}{} dan
\mathl{\varphi(P) \in V_2 \subset \R^n}{,}
beserta suatu perluasan
\maabbdisp {\tilde{\varphi}} {V_1 } { V_2
} {}
yang merupakan difeomorfisme.

}
{} {}

\inputaufgabe
{}
{

Misalkan cakram tertutup
\mathl{B  \left(  0 , 1 \right)}{} dilengkapi dengan
\definitionsverweis {orientasi standar}{}{}
dari $\R^2$. Apakah orientasi yang ditentukan oleh
\definitionsverweis {normal luar}{}{}
pada batas
\zusatzklammer {yaitu pada lingkaran satuan} {} {}
searah atau berlawanan arah jarum jam?

}
{} {}

  \zwischenueberschrift{Soal untuk dikumpulkan}

\bild{ \begin{center}
\includegraphics[width=5.5cm]{ \bildeinlesungsvg {Not-star-shaped} {svg} }
\end{center}
\bildtext {} }

\bildlizenz { Not-star-shaped.svg } {} {WillT.Net} {Commons} {PD} {}

\inputaufgabe
{4}
{

Untuk anulus

\mavergleichskettedisp
{\vergleichskette
{M
}
{ =} { \left\{  x \in \R^2  \mid   1 \leq \Vert { x } \Vert \leq 2         \right\}
}
{ } {
}
{ } {
}
{ } {
}
}
{}{}{}
berikan secara eksplisit
\definitionsverweis {peta}{}{}
yang menunjukkan bahwa $M$ merupakan
\definitionsverweis {manifold berbatas}{}{.}

}
{} {}

\inputaufgabe
{6}
{

Tunjukkan bahwa
\definitionsverweis {setengah bidang}{}{} $\R_{\geq 0} \times \R$ dan
\definitionsverweis {kuadran}{}{}
\mathl{\R_{\geq 0} \times \R_{\geq 0}}{} tidak
\definitionsverweis {difeomorfik}{}{.}

}
{(Apa pengertian difeomorfisme yang sesuai dalam konteks ini?)} {}

\inputaufgabe
{6}
{

Misalkan
\mathl{M= B  \left(  0 , 1 \right) \setminus \{ (0,1), (0,-1)\} \subseteq \R^2}{,}
yaitu
\definitionsverweis {cakram tertutup}{}{}
yang dua titik batasnya telah dibuang. Misalkan
\mathl{N={]{-1},1[} \times  [-1,1]}{} merupakan produk sebuah interval terbuka dan sebuah interval tertutup. Tunjukkan bahwa
\mathkor {} {M} {dan} {N} {}
merupakan
\definitionsverweis {manifold berbatas}{}{}
yang saling
\definitionsverweis {difeomorfik}{}{.}

}
{} {}

\inputaufgabe
{4}
{

Misalkan
\mathl{V_1,V_2 \subseteq \R^n}{}
merupakan
\definitionsverweis {himpunan bagian terbuka}{}{}
dan
\maabbdisp {\varphi} {V_1 } { V_2
} {}
suatu
\definitionsverweis {difeomorfisme}{}{}
yang menginduksi
\definitionsverweis {homeomorfisme}{}{}
antara
\mathkor {} {V_1 \cap H} {dan} {V_2 \cap H} {}
dan, dengan demikian, juga antara
\mathkor {} {V_1 \cap \partial H} {dan} {V_2 \cap \partial H} {}
\zusatzklammer {$H$ menyatakan
\definitionsverweis {setengah ruang}{}{}
dan $\partial H$ batasnya} {} {.}
Tunjukkan bahwa pembatasan pada batas juga merupakan
\definitionsverweis {difeomorfisme}{}{.}

}
{} {}

\inputaufgabe
{4}
{

Misalkan bola satuan tertutup
\mathl{B  \left(  0 , 1 \right)}{} dilengkapi dengan
\definitionsverweis {orientasi standar}{}{} dari $\R^3$. Tentukan apakah kedua vektor singgung
\mathkor {} {(2,1,0)} {dan} {(3,-1,0)} {}
di kutub utara
\mathl{(0,0,1)}{} merepresentasikan orientasi pada batas yang diinduksi oleh
\definitionsverweis {normal luar}{}{}
\zusatzklammer {yaitu pada permukaan bola satuan} {} {,}
atau tidak.

}
{} {}

 [[Kategorie:Latexseite]]
